% -----------------------------------------------------------------------
\section{Preliminaries}
\label{sec:prelim}
% -----------------------------------------------------------------------

This section fixes the problem setting: the flow-state and geometry
representation, the stochastic model of the inflow forcing, and the ensemble
data-assimilation scheme that the neural surrogate is later embedded in.

\subsection{Domain and State Representation}
\label{sec:domain}

The computational domain covers an $896\,\mathrm{m} \times 896\,\mathrm{m}
\times 128\,\mathrm{m}$ section of the Barcelona urban block, discretised onto a
uniform Cartesian grid of $224 \times 224 \times 32$ cells at $4\,\mathrm{m}$
resolution. At this resolution individual building facades are represented by
two-to-four cell widths, capturing the dominant flow-blocking and street-canyon
channelling effects while keeping the state dimensionality manageable.

The flow state is the three-component velocity field
\begin{equation}
  \state \in \Real^{3 \times 32 \times 224 \times 224},
  \label{eq:state}
\end{equation}
comprising horizontal components $u$ (west--east) and $v$ (south--north) and
the vertical component $w$. Building geometry is encoded by a binary mask
$\mask \in \{0,1\}^{32 \times 224 \times 224}$, where $g_{ijk}=1$ denotes
a fluid cell and $g_{ijk}=0$ denotes an obstacle cell. The mask is fixed for
a given domain and does not vary between ensemble members or time steps.

The inflow forcing is characterised by two parameters
$\bth = (\iang,\,\ivel)$,
where $\iang$ is the inflow wind angle and $\ivel$ is the magnitude of the inflow
velocity. These
parameters vary across ensemble members and, within a member, evolve over time
according to an AR(2) process (Section~\ref{sec:prior}).

\subsection{Inflow Parameterisation}
\label{sec:prior}

The inflow parameters evolve in time according to an AR(2) relaxation process:
\begin{equation}
  \bth_{k+1}
    = \mu_{\mathrm{ext}}
      + \phi_1(\bth_k - \mu_{\mathrm{ext}})
      + \phi_2(\bth_{k-1} - \mu_{\mathrm{ext}})
      + \sigma_\eta\,\boldsymbol{\eta}_k,
  \label{eq:ar2}
\end{equation}
where $\mu_{\mathrm{ext}}$ is an external climatological mean, $\phi_1$ and
$\phi_2$ are the autoregressive coefficients that control the correlation
length of the parameter trajectory, and $\sigma_\eta\boldsymbol{\eta}_k$ is
white process noise. The stationary distribution of the AR(2) process provides
the marginal prior on $\bth$ at any given time step. For the
inflow angle the prior marginal is $\iang \sim \mathcal{N}(0,\,6^{\circ\,2})$,
reflecting the typical range of wind direction variability in the Barcelona
metropolitan area. For the velocity magnitude the prior is
$\ivel \sim \mathcal{N}(7.5\,\mathrm{m/s},\,0.5^2\,\mathrm{m}^2/\mathrm{s}^2)$,
consistent with moderate urban boundary-layer winds. The AR(2) structure ensures
temporal coherence of the parameter trajectories within each ensemble member,
preventing unphysical step-changes in inflow conditions between successive time
steps.

\subsection{Ensemble Smoother with Multiple Data Assimilation}
\label{sec:esmda_framework}

The Ensemble Smoother with Multiple Data Assimilation (\esmda{}) of
\cite{emerick2013} is an iterative ensemble Kalman update scheme that performs
$\niter$ analysis passes per assimilation window, each pass inflating
the observation error covariance by a factor $\alpha$ chosen so that
$\sum_{i=1}^{\niter} 1/\alpha_i = 1$. The ensemble update at each
pass is
\begin{equation}
  \xctrl^{a,n}
    = \xctrl^{f,n}
      + \Kgain
        \bigl(\yobs + \sqrt{\alpha}\,\boldsymbol{\varepsilon}^n
              - \Hobs\xctrl^{f,n}\bigr),
  \label{eq:esmda}
\end{equation}
where $\xctrl^{f,n}$ and $\xctrl^{a,n}$ are the forecast and analysis
augmented state vectors for member $n$, $\yobs$ is the observation vector,
$\Hobs$ is the linear observation operator, $\boldsymbol{\varepsilon}^n
\sim \mathcal{N}(\mathbf{0},\Rcov)$ is perturbed noise, and the Kalman
gain is
\begin{equation}
  \Kgain
    = \mathbf{C}_{xH}
      \bigl(\mathbf{C}_{HH} + \alpha\Rcov\bigr)^{-1},
  \label{eq:gain}
\end{equation}
with $\mathbf{C}_{xH}$ the ensemble cross-covariance between the augmented
state and the mapped state $\Hobs\xctrl^f$, and $\mathbf{C}_{HH}$
the ensemble covariance of the mapped state. The augmented control vector
$\xctrl$ can in general concatenate the velocity field and the inflow
parameter vector for a joint state--parameter update.
